\documentclass[12pt]{article}
\usepackage[utf8]{inputenc}
\usepackage{amsmath}
\usepackage[margin=0.5in]{geometry}
\usepackage{multicol}

\pagestyle{empty}

\begin{document}


Nombre y Apellido \rule{10cm}{0.4pt}

\begin{enumerate}
    \item (3 puntos) Un sistema termodinámico recibe 500 J de calor y realiza un trabajo de 200 J sobre su entorno. ¿Cuál es el cambio en la energía interna del sistema en kcal?

    \item (2 puntos) Un gas se expande y realiza un trabajo de 150 J mientras absorbe 600 J de calor. Calcula el cambio en la energía interna del gas en kcal.

  \item (3 puntos) Durante una compresión adiabática, un gas realiza un trabajo de 250 J y su energía interna aumenta en 250 J. ¿Cuánto calor ha intercambiado el gas con el entorno? Justifica tu respuesta.

\end{enumerate}

Nombre y Apellido \rule{10cm}{0.4pt}

\begin{enumerate}

    \item (3 puntos) Un sistema realiza un trabajo de 300 J y su energía interna disminuye en 100 J. ¿Cuánto calor ha intercambiado el sistema con su entorno en kgm?

    \item (3 puntos) Un gas se expande y realiza un trabajo de 150 J mientras absorbe 600 J de calor. Calcula el cambio en la energía interna del gas en kgm.

    \item (2 puntos) Durante una compresión adiabática, un gas realiza un trabajo de 250 J y su energía interna aumenta en 250 J. ¿Cuánto calor ha intercambiado el gas con el entorno? Justifica tu respuesta.

\end{enumerate}




\end{document}
